%-----------------------------------------------------------------------
% Template File for Science China Information Sciences
% Downloaded from http://scis.scichina.com
% Please compile the tex file using LATEX or PDF-LATEX or CCT-LATEX
%-----------------------------------------------------------------------

\documentclass{SCIS2026}
%%%%%%%%%%%%%%%%%%%%%%%%%%%%%%%%%%%%%%%%%%%%%%%%%%%%%%%
%%% Author's definitions for this manuscript
%%% 作者附加的定义
%%% 常用环境已经加载好, 不需要重复加载
%%%%%%%%%%%%%%%%%%%%%%%%%%%%%%%%%%%%%%%%%%%%%%%%%%%%%%%
\usepackage{threeparttable}
\usepackage{colortbl}  %彩色表格需要加载的宏包
\usepackage{xcolor}
\usepackage{array}   %对表列和表格线的设置需要用到array宏包
\usepackage{algorithm}
\usepackage{algorithmic}
\usepackage{bm}
\usepackage{amsmath}
%\usepackage{subcaption}
%%%%%%%%%%%%%%%%%%%%%%%%%%%%%%%%%%%%%%%%%%%%%%%%%%%%%%%
%%% Begin. 开始
%%%%%%%%%%%%%%%%%%%%%%%%%%%%%%%%%%%%%%%%%%%%%%%%%%%%%%%
\begin{document}
%\oa
%%%%%%%%%%%%%%%%%%%%%%%%%%%%%%%%%%%%%%%%%%%%%%%%%%%%%%%
%%% Authors do not modify the information below
%%% 作者不需要修改此处信息
\ArticleType{RESEARCH PAPER}
%\SpecialTopic{}
\Year{2025}
\Month{January}
\Vol{68}
\No{1}
\DOI{}
\ArtNo{}
\ReceiveDate{}
\ReviseDate{}
\AcceptDate{}
\OnlineDate{}
\AuthorMark{}
\AuthorCitation{}
%%%%%%%%%%%%%%%%%%%%%%%%%%%%%%%%%%%%%%%%%%%%%%%%%%%%%%%

%%% title: 标题
%%%   \title{title}{title for citation}
\title{Argus: Towards Reliable Network Intrusion Detection in Noisy and Open Environments}{Title for citation}

%%% responding author: 通信作者
%%%   \author[number]{Full name}{{email@xxx.com}}
%%% General author: 一般作者
%%%   \author[number]{Full name}{}
%%% Equal Contribution: 同等贡献作者
%%%   \author[number\dag]{Full name}{}
\author[1,2]{Qianwei Meng}{}
\author[1,2]{Qingjun Yuan}{{gcxyuan@outlook.com}}
\author[1,2]{Guangsong Li}{}
\author[1,2]{Yongjuan Wang}{}

%%% Authors' contribution. 同等贡献声明
%\contributions{These authors contributed equally to this work.}

%%% Address. 地址
%%%   \address[number]{Affiliation, City Postcode, Country}
\address[1]{Information Engineering University, Zhengzhou 450001, China}
\address[2]{Henan Key Laboratory of Network Cryptography Technology, Zhengzhou 450001, China}
\address[3]{Affiliation, City 000000, Country}

%%% Abstract. 摘要
\abstract{An abstract (100--200 words) is a summary of the content of the manuscript. It should briefly describe the research purpose, method, result and conclusion. The extremely professional terms, special signals, figures, tables, chemical structural formula, and equations should be avoided here, and citation of references is not allowed. An abstract (100--200 words) is a summary of the content of the manuscript. It should briefly describe the research purpose, method, result and conclusion. The extremely professional terms, special signals, figures, tables, chemical structural formula, and equations should be avoided here, and citation of references is not allowed. An abstract (100--200 words) is a summary of the content of the manuscript. It should briefly describe the research purpose, method, result and conclusion. The extremely professional terms, special signals, figures, tables, chemical structural formula, and equations should be avoided here, and citation of references is not allowed. An abstract (100--200 words) is a summary of the content of the manuscript. It should briefly describe the research purpose, method, result and conclusion. The extremely professional terms, special signals, figures, tables, chemical structural formula, and equations should be avoided here, and citation of references is not allowed.}

%%% Keywords. 关键词
\keywords{Network intrusion detection system, unknown attack, encrypted traffic}

\maketitle


%%%%%%%%%%%%%%%%%%%%%%%%%%%%%%%%%%%%%%%%%%%%%%%%%%%%%%%
%%% The main text. 正文部分
%%%%%%%%%%%%%%%%%%%%%%%%%%%%%%%%%%%%%%%%%%%%%%%%%%%%%%%
\section{Introduction}\label{Introduction}
%Please use this sample as a guide for preparing your article. Please read all of the following manuscript preparation instructions carefully and in their entirety. The manuscript must be in good scientific American English; this is the author's responsibility. All files will be submitted through our online electronic submission system at \href{https://mc03.manuscriptcentral.com/scis}{HERE}.

The effectiveness of Network Intrusion Detection Systems (NIDS) highly relies on large-scale and high-quality annotated traffic datasets. However, obtaining accurate and stable annotation knowledge in real-world complex network environments is facing unprecedented challenges \cite{guerra2022datasets}. As attackers widely adopt techniques such as traffic shaping, protocol obfuscation, and heavily encrypted payloads, malicious and normal traffic highly overlap in statistical dimensions and behavioral patterns \cite{EC-Net}. Such deep semantic confusion makes it extremely easy for both domain experts and automated rule engines to deviate during the annotation process, thereby introducing \textbf{label noise} \cite{kim2022human, FlowRefiner}. In such scenarios, label noise severely distorts the decision boundaries of known categories, interferes with the model's feature extraction of the underlying traffic structure, and substantially degrades the recognition accuracy of NIDS.

Beyond internal annotation distortion, NIDS in open environments face the severe challenge of unknown attack traffic \cite{SEEN}. These unknown threats typically originate from the outbreak of zero-day vulnerabilities or the mutated evolution of existing attack patterns. Since these sample categories are unseen during training, they manifest as out-of-distribution anomalies during the model inference stage. More severely, there is a significant compound negative effect between label noise and unknown attacks: the feature space shift caused by label noise suppresses the model's generalization upper bound, causing it to produce ``over-confident" misclassifications when confronting unknown traffic \cite{add_3_2, add_3_3}. Experimental evidence shows that higher noise ratios lead to worse NIDS performance. During dynamic system updates, if these misdetected samples are reinjected as ``pseudo-labels," a vicious feedback loop of ``label noise - unknown misdetection - secondary pollution" will be triggered. This coupling effect not only leads to the accelerated degradation of detection performance but also fundamentally undermines the reliability of NIDS in long-term deployment tasks.


To address label noise, existing methods are mainly categorized into loss correction and sample selection. Loss correction attempts to compensate for the interference by modifying the loss function or estimating the transition matrix. However, these methods are mostly based on simple ``instance-independent" assumptions, making it difficult to handle complex ``instance-dependent" noise in real traffic, and they are prone to overfitting in high-noise environments. Sample selection methods utilize the ``small-loss" criterion to filter data; yet, in realistic asymmetric noise scenarios, the loss distributions of clean and noisy samples highly overlap. This causes the model to easily incorrectly include high-confidence noisy samples while discarding critical ``hard clean samples," ultimately leading to ambiguous decision boundaries and weakened discriminative performance. Regarding unknown attacks, existing methods carefully design various open-set detection models or extreme value analysis classification heads, such as OpenMax. Nevertheless, these methods generally assume that the data labels used for initial training are absolutely clean. Once deployed in real network environments characterized by both label pollution and open-set unknown traffic, their relied deep feature spaces are already contaminated. Consequently, the corresponding detection boundaries and confidence measurement mechanisms easily collapse, completely failing to break the aforementioned vicious cycle.

In this paper, we find that introducing complementary label learning combined with a confidence screening mechanism in the early stage of model training can prevent the model from forcefully fitting incorrect annotations, thereby effectively solving the problems of overfitting and feature space pollution caused by label noise. Furthermore, breaking through the traditional discriminative paradigm based on Softmax maximum confidence, we design a multidimensional principal subspace detection method for unknown attack traffic. By characterizing the deviation between deep features and known traffic distributions from a geometric perspective, this method can highly stably resolve the rejection and detection of unknown attacks.

Specifically, we propose \textbf{Argus}, a novel label denoising model for unknown encrypted traffic detection. The overall model consists of two phases. Phase 1 is the robust feature learning phase under noisy labels, comprising three progressive sub-stages: complementary label learning, selective complementary training, and high-confidence sample fitting. This phase guides the model to progressively purify the training data and extract robust representations without relying on accurate prior annotations. Phase 2 is the unknown encrypted traffic detection phase, which utilizes the robust deep features obtained in Phase 1 to construct multidimensional principal singular value subspaces for each known category. Through an original multi-subspace angular deviation, scores are calculated for test traffic, and any sample whose score exceeds the threshold-estimated boundary is identified as unknown traffic. Extensive experiments on multiple encrypted traffic and attack datasets demonstrate the significant superiority of our model in dealing with varied asymmetric and symmetric noise, as well as open-network unknown traffic.


The main contributions of this paper are summarized as follows:

\begin{itemize}
	\item We propose Argus, a noise-robust learning framework for encrypted traffic. By introducing a joint training strategy based on complementary labels, we effectively prevent the decision boundary distortion caused by the model's premature fitting of noisy labels, without relying on complex noise distribution assumptions.
	
	\item We design an unknown attack detection mechanism based on multi-subspace angular deviation. Different from traditional single-center or confidence-based methods, we utilize the geometric projection deviation within the feature space to effectively suppress the ``over-confidence" issue inherent in neural networks, maintaining a clear known/unknown boundary even in high-noise scenarios.
	
	\item We empirically validate the effectiveness of Argus across multiple public datasets. Compared with existing state-of-the-art (SOTA) methods, the proposed model not only maintains remarkably high classification accuracy for known categories under extreme noise environments of up to 50\%, but also exhibits outstanding open-world unknown attack detection capabilities and great generalizability.
	
\end{itemize}

The remainder of this paper is structured as follows: 
Section~\ref{Related_work} reviews related studies on noisy-label handling and unknown malicious traffic detection. 
Section~\ref{Method} presents the proposed Argus framework. 
Section~\ref{Experiment} reports the experimental results and analysis. 
Finally, Section~\ref{conclusion} concludes this paper.



\section{Related work}\label{Related_work}

Please use 
\subsection{Noise label Handling}

\cite{FlowRefiner}
\cite{BAPTISM}


\subsection{Unknown Malicious Traffic Detection}
\cite{8_HyperVision}

\cite{FEC-OSL}

\cite{liu2023good}

\cite{luoxiang}

\subsection{In summary}
Existing research has made significant progress in addressing noisy labels and unknown malicious traffic detection. 



\section{Method}\label{Method}
As shown in Figure \ref{Argus}, Argus is a two-phase framework for robust known encrypted traffic classification and reliable unknown traffic detection under noisy-label conditions. Phase 1 performs noise-robust feature learning through the SelCLL-PL framework, which progressively integrates Complementary Label Learning (CLL), Selective Complementary Label Learning (SelCLL), and Selective Positive Learning (SelPL). Specifically, CLL first introduces exclusion-based supervision to avoid early fitting to incorrect labels; SelCLL then selects relatively reliable samples based on prediction confidence; and SelPL further fits positive labels on high-confidence samples to purify training and learn discriminative representations of known traffic. Phase 2 targets open-world unknown traffic detection. It extracts robust deep features from the Phase 1 model, constructs a multidimensional principal subspace for each known class via SVD, and uses the maximum projection similarity between a test sample and all known-class subspaces to compute a multi-subspace angular deviation score for unknown detection.

%As shown in Figure \ref{Argus}, the proposed Argus model consists of two phases, aiming to achieve robust classification of known encrypted traffic and reliable detection of unknown encrypted traffic under noisy-label conditions. Phase 1 focuses on robust feature learning with noisy labels and comprises three progressive sub-stages: Complementary Label Learning (CLL), Selective Complementary Label Learning (SelCLL), and Selective Positive Learning (SelPL), which together form the SelCLL-PL framework. In this phase, exclusion-based complementary supervision is first introduced to prevent the model from prematurely fitting incorrect labels. Confidence-based sample selection is then employed to progressively identify reliable samples, and positive label fitting is further performed on high-confidence samples, thereby purifying the training process and learning stable representations for known traffic classes. Phase 2 is designed for unknown encrypted traffic detection in open environments. Specifically, Argus extracts deep features of known-class traffic using the robust model trained in Phase 1 and constructs a multidimensional principal subspace for each known class via SVD. During testing, the model computes the maximum projection similarity between an input sample and all known-class subspaces, based on which a multi-subspace angular deviation score is defined as the unknown detection score.


\begin{figure*}[ht]
	\centering
	\includegraphics[width=0.98\textwidth]{img/1}
	\setlength{\abovecaptionskip}{0cm}
	\caption{The structure diagram of Argus.}
	\label{Argus}
	\vspace{-0.6cm}       
\end{figure*}


\subsection{Phase 1: Robust Training under Noisy Labels}
Phase 1 aims to achieve robust training of the encrypted traffic classification model under noisy-label conditions. As illustrated in Figure \ref{Argus}, this phase adopts a three-stage training strategy. 

%First, Stage 1 employs Complementary Label Learning (CLL) to prevent the model from directly fitting potentially incorrect observed labels in the early training stage. It further applies a basic confidence threshold to select relatively reliable samples for continued complementary-label training. Subsequently, Stage 2 introduces positive label fitting on high-confidence samples, enabling the model to fully exploit reliable supervisory information. Finally, Stage 3 leverages the confidence separation capability developed after Stage 2 to explicitly partition the original training set, and performs robust retraining using only high-confidence clean samples. Through this progressive process, the model can gradually reduce the interference of noisy labels during parameter optimization and learn more stable and discriminative representations for encrypted traffic.


\subsubsection{Stage 1}
Stage 1 consists of two consecutive steps. First, complementary label learning is performed on all noisy training samples. Then, sample selection is conducted based on the model's prediction confidence for the observed labels, where only samples whose confidence is higher than the random-guessing level are retained for further complementary-label training. The objective of this stage is to learn preliminary robust traffic representations without directly relying on the correctness of the observed labels, thereby laying the foundation for subsequent high-confidence sample fitting.

\paragraph{Traffic Representation Learning Based on Complementary Label Learning}

Conventional encrypted traffic classification models typically rely on standard supervised learning, where each traffic sample is trained to match its observed label. While effective with clean annotations, this strategy becomes vulnerable in real-world networks where labels may be noisy. Directly fitting incorrect labels can inject misleading supervision into the feature space and lead to overfitting.

To address this issue, we adopt Complementary Label Learning \cite{CLL-1, CLL-2, CLL-3}. Instead of enforcing a sample to belong to its observed class, CLL provides exclusion-based supervision by teaching the model that the sample does not belong to a selected complementary class, thereby reducing the risk of fitting noisy labels.

Assume that the encrypted traffic classification task contains $c$ known classes, and the label space is defined as $\mathcal{Y}=\{1,2,\ldots,c\}$. Given a noisy training set $\mathcal{D}_{tr}=\{(x_i,\tilde{y}_i)\}_{i=1}^{N}$, where $x_i$ denotes the $i$-th sample, and $\tilde{y}_i \in \mathcal{Y}$ denotes its observed label. Due to weak annotation errors or category confusion in realistic network environments, $\tilde{y}_i$ is not necessarily reliable.

Let the classification model be denoted as $f(x;\theta)$, where $\theta$ represents the model parameters. The model maps an input sample into a $c$-dimensional class-score space and then obtains the class probability distribution through the softmax function:
\begin{equation}
	p_i=\mathrm{softmax}(f(x_i;\theta)), \quad p_i \in \mathbb{R}^{c},
\end{equation}
where $p_{i,k}$ denotes the probability that sample $x_i$ is predicted as class $k$, satisfying $\sum_{k=1}^{c} p_{i,k}=1$.

In conventional positive label fitting, the cross-entropy loss is commonly adopted:
\begin{equation}
	\mathcal{L}_{CE}(f,\tilde{y}_i)
	=
	-\sum_{k=1}^{c}\mathbb{I}(\tilde{y}_i=k)\log p_{i,k},
	\label{eq:CE_loss}
\end{equation}
where $\mathbb{I}(\cdot)$ is the indicator function. This loss encourages the model to increase the predicted probability of the observed class $\tilde{y}_i$, namely, $p_{i,\tilde{y}_i} \rightarrow 1$.

When the observed label is correct, this optimization direction can effectively enhance class discrimination. However, when the observed label is noisy, the model is forced to pull the sample feature toward an incorrect class. CLL adopts a different training objective. For a sample $(x_i,\tilde{y}_i)$, a complementary label is first randomly sampled from the label set excluding the observed label:
\begin{equation}
	\bar{y}_i \sim \mathrm{Uniform}(\mathcal{Y}\setminus\{\tilde{y}_i\}).
\end{equation}
The model is then trained to reduce the predicted probability of this complementary class, i.e., $p_{i,\bar{y}_i} \rightarrow 0$.

The corresponding complementary-label loss is defined as
\begin{equation}
	\mathcal{L}_{CLL}(f,\bar{y}_i)
	=
	-\sum_{k=1}^{c}\mathbb{I}(\bar{y}_i=k)\log(1-p_{i,k}).
	\label{eq:CLL_loss}
\end{equation}

Since $\bar{y}_i$ is randomly selected from classes other than the observed label, even if $\tilde{y}_i$ is incorrect, $\bar{y}_i$ still has a high probability of not being the true class of the sample. In other words, compared with directly fitting the observed label, complementary label learning is less likely to inject erroneous target-class supervision into the model.

%For example, suppose an encrypted flow that truly belongs to online gaming traffic is incorrectly labeled as file transfer traffic. Conventional positive label fitting would directly push the model to recognize this sample as file transfer traffic. In contrast, complementary label learning may randomly sample ``video traffic'' as the complementary label and train the model to regard this sample as not belonging to video traffic. Although the original observed label is incorrect, the supervisory signal that ``this sample is not video traffic'' is still highly likely to be correct. Therefore, complementary label learning can effectively reduce the misleading effect of noisy labels and improve the robustness of encrypted traffic classification models under erroneous annotation scenarios.

\paragraph{Confidence-Based Selective Complementary Label Training}
Although CLL alleviates overfitting to noisy labels, its exclusion-based supervision is indirect and usually converges more slowly than positive label fitting. To improve training efficiency and reduce the influence of unreliable samples, we introduce a confidence-based selection mechanism after the initial CLL stage. Specifically, the current model's confidence on the observed label is used to estimate sample reliability, and only relatively reliable samples are retained for subsequent complementary-label training.

For a $c$-class classification task, a model without class preference assigns an average probability of approximately $1/c$ to each class. We therefore use $1/c$ as the basic confidence threshold and define the selected sample set as
\begin{equation} 
	\mathcal{D}_{CLL}^{sel} 
	=
	\left\{ 
	(x_i,\tilde{y}_i)\in \mathcal{D}_{tr} 
	\mid
	p_{i,\tilde{y}_i}>\frac{1}{c} 
	\right\}. 
	\label{eq:selected_cl_set} 
\end{equation} 

A sample satisfying $p_{i,\tilde{y}_i}>1/c$ receives higher-than-random confidence on its observed label, indicating stronger label consistency and higher reliability. In contrast, samples with $p_{i,\tilde{y}_i}\leq 1/c$ are likely to be mislabeled, ambiguous, or insufficiently learned, and are temporarily excluded from parameter updates.

The model is then optimized with the same complementary-label loss, but only on the selected subset:
\begin{equation} 
	\min_{\theta} 
	\frac{1}{|\mathcal{D}_{CLL}^{sel}|} 
	\sum_{(x_i,\tilde{y}_i)\in \mathcal{D}_{CLL}^{sel}} 
	\mathcal{L}_{CLL}(f,\bar{y}_i). 
	\label{eq:selected_cl_objective} 
\end{equation} 

As shown in Figure \ref{pro}(a), after initial CLL training, clean and noisy samples start to show different confidence distributions on their observed labels. The threshold $1/c$ filters out low-confidence samples, enabling subsequent CLL training to rely on more reliable data while preserving robustness. Consequently, as illustrated in Figure \ref{pro}(b), Stage 1 produces a clearer separation between clean and noisy samples, improving convergence stability and noise suppression.


\begin{figure*}[ht]
	\centering
	\includegraphics[width=1.0\textwidth]{img/pro}
	\setlength{\abovecaptionskip}{0cm}
	\caption{Prediction confidence distributions of clean and noisy samples under the Sym-10\% setting on the MAL\_TLS2023 dataset, where orange denotes noisy samples and blue denotes clean samples.}
	\label{pro}
	\vspace{-0.8cm}
\end{figure*}


\subsubsection{Stage 2}

After Stage 1, the model has obtained preliminary noise-suppression capability and can partially separate clean and noisy samples according to the confidence assigned to their observed labels. However, CLL only provides exclusion-based supervision and lacks direct positive constraints on the target class. For samples with reliable labels, positive label fitting remains more effective and efficient. Therefore, Stage 2 introduces Selective Positive Learning, which performs direct supervised training only on high-confidence samples.

The key assumption is that, after CLL and SelCLL, label-consistent samples tend to receive higher probabilities on their observed labels, while mislabeled or semantically inconsistent samples usually have lower confidence. Thus, we use a stricter confidence threshold $\gamma$ to select reliable samples:
\begin{equation} 
	\mathcal{D}_{CE}^{sel} 
	=
	\left\{ 
	(x_i,\tilde{y}_i)\in \mathcal{D}_{tr} 
	\mid
	p_{i,\tilde{y}_i}>\gamma
	\right\}, 
	\label{eq:stage2_selected_set} 
\end{equation} 
where $\gamma$ denotes the high-confidence threshold. In this paper, we set $\gamma=0.1$, and only samples whose observed-label confidence exceeds this threshold are used for positive label fitting.

In this stage, the optimization objective changes from suppressing the complementary class to increasing the probability of the observed target class. The model is trained with the cross-entropy loss on the selected subset:
\begin{equation} 
	\min_{\theta} 
	\frac{1}{|\mathcal{D}_{CE}^{sel}|} 
	\sum_{(x_i,\tilde{y}_i)\in \mathcal{D}_{CE}^{sel}} 
	\mathcal{L}_{CE}(f,\tilde{y}_i). 
	\label{eq:stage2_ce_objective} 
\end{equation} 

Since the selected samples have higher label confidence, the risk of directly fitting noisy labels is greatly reduced. As shown in Figure \ref{pro}(c), the threshold $\gamma$ further filters out low-confidence noisy samples while retaining high-confidence clean samples for positive supervision. This process enlarges the confidence gap between clean and noisy samples, enabling the model to better exploit reliable annotations and improve class discriminability while maintaining robustness to label noise.


\subsubsection{Stage 3}
After Stage 2, positive label fitting on high-confidence samples further enlarges the confidence gap between clean and noisy samples. As shown in Figure \ref{pro}(c), label-consistent samples usually obtain high prediction confidence on their observed labels, whereas mislabeled samples tend to have low confidence due to the semantic inconsistency between traffic features and annotations. Therefore, we use the model obtained after Stage 2 to explicitly partition the original noisy training set, and retain only high-confidence samples for final supervised training.

Let $\theta^{(2)}$ denote the model parameters after Stage 2. For each sample $(x_i,\tilde{y}_i)\in \mathcal{D}_{tr}$, the predicted probability distribution is computed as
\begin{equation} 
	p_i^{(2)} 
	=
	\mathrm{softmax}\left(f(x_i;\theta^{(2)})\right).
	\label{eq:stage3_prob}
\end{equation} 
The confidence of its observed label is denoted by $p_{i,\tilde{y}_i}^{(2)}$.

Using the same high-confidence threshold $\gamma$ as in Stage 2, the original training set is divided into a clean set and a noisy set:
\begin{equation} 
	\mathcal{D}_{clean} 
	=
	\left\{ 
	(x_i,\tilde{y}_i)\in \mathcal{D}_{tr} 
	\mid
	p_{i,\tilde{y}_i}^{(2)}>\gamma
	\right\},
	\label{eq:clean_set}
\end{equation} 
\begin{equation} 
	\mathcal{D}_{noisy} 
	=
	\left\{ 
	(x_i,\tilde{y}_i)\in \mathcal{D}_{tr} 
	\mid
	p_{i,\tilde{y}_i}^{(2)}\leq \gamma
	\right\}.
	\label{eq:noisy_set}
\end{equation} 
Here, $\mathcal{D}_{clean}$ contains samples regarded as reliable, while $\mathcal{D}_{noisy}$ contains low-confidence samples that are excluded from subsequent optimization. Different from Stage 2, which selects samples during training, this stage explicitly filters the original training set after convergence, thereby reducing the interference of suspected noisy labels in the final model refinement.

The model is then initialized with $\theta^{(2)}$ and further trained only on $\mathcal{D}_{clean}$ using the cross-entropy loss:
\begin{equation} 
	\theta^{(3)} 
	=
	\arg\min_{\theta} 
	\frac{1}{|\mathcal{D}_{clean}|} 
	\sum_{(x_i,\tilde{y}_i)\in \mathcal{D}_{clean}} 
	\mathcal{L}_{CE}(f,\tilde{y}_i),
	\quad
	\theta \leftarrow \theta^{(2)}.
	\label{eq:stage3_objective}
\end{equation} 

The resulting parameters $\theta^{(3)}$ form the final robust classification model in Phase 1. Since the retained samples have high observed-label confidence, the model can refine its decision boundaries on a cleaner data distribution and learn more stable traffic representations instead of memorizing noisy correlations. As shown in Figure \ref{pro}(d), after Stage 3, clean samples are concentrated near high confidence, while most noisy samples remain in the low-confidence region.

This final refinement further improves robustness to noisy labels and provides more reliable feature representations for class subspace construction and unknown traffic detection in Phase 2. 
%The overall training process is summarized in Algorithm \ref{alg:phase1}.




%\begin{algorithm}
%	\footnotesize
%	\caption{Training Process of the SelCLL-PL Framework (Phase 1)}
%	\label{alg:phase1}
%	\begin{algorithmic}[1]
%		\REQUIRE Noisy training set $D_{tr}$; class number $c$; threshold $\gamma$; epochs $N_1,N_2,N_3$.
%		\ENSURE Robust classifier $F=h\circ E$.
%		
%		\STATE Initialize model parameters $\theta$.
%		
%		\STATE \textbf{/* Stage 1: CLL + SelCLL */}
%		\FOR{$epoch=1$ to $N_1$}
%		\STATE Select training samples by $p_i(\tilde{y}_i)>1/c$ after warm-up; otherwise use all samples.
%		\STATE Sample complementary labels from $\{1,\ldots,C\}\setminus\{\tilde{y}_i\}$.
%		\STATE Optimize $\theta$ using complementary label loss $L_{\text{CLL}}$.
%		\ENDFOR
%		
%		\STATE \textbf{/* Stage 2: SelPL */}
%		\FOR{$epoch=1$ to $N_2$}
%		\STATE Select high-confidence samples with $p_i(\tilde{y}_i)\geq\gamma$.
%		\STATE Optimize $\theta$ using cross-entropy loss $L_{\text{CE}}$ on selected samples.
%		\ENDFOR
%		
%		\STATE \textbf{/* Stage 3: Robust Re-training */}
%		\STATE Partition $\tilde{V}$ into clean set $\mathcal{V}_{\text{clean}}$ and noisy set by threshold $\gamma$.
%		\FOR{$epoch=1$ to $N_3$}
%		\STATE Re-train the model on $\mathcal{V}_{\text{clean}}$ using $L_{\text{CE}}$.
%		\ENDFOR
%		
%		\STATE \textbf{return} $F=h\circ g$.
%	\end{algorithmic}
%\end{algorithm} 



\subsection{Phase 2: Unknown Traffic Detection via Class-Wise Multidimensional Subspaces}

To enhance the rejection capability for unknown-class traffic in open environments, we further introduce a class-wise multidimensional subspace-based unknown traffic detection method after Phase 1. 

%This method does not modify the training objective of the classification model in Phase 1, but serves as a feature-space detection module after training. 

%Specifically, the robust model obtained in Phase 1 is first used to extract deep features of known-class traffic. Then, based on the feature representations of high-confidence known samples, a multidimensional principal subspace is constructed for each known class. Finally, during testing, the geometric consistency between the test sample feature and all known-class subspaces is computed, and unknown traffic is identified according to its angular deviation.


\subsubsection{Build Class Subspaces from the Phase 1 Model}

Let the known classes observed during training be $\mathcal{Y}_{in}=\{1,2,\ldots,c\}$, where $c$ is the number of known classes. The noisy training set is denoted as $\mathcal{D}_{tr}=\{(x_i,\tilde{y}_i)\}_{i=1}^{N}$, with $\tilde{y}_i\in\mathcal{Y}_{in}$. After Phase 1, we obtain the final robust model parameters $\theta^{(3)}$.

The trained network is decomposed into a feature extractor and a classifier, i.e., $F(x)=h_{\phi}(g_{\theta}(x))$, where $g_{\theta}(\cdot)$ extracts deep features and $h_{\phi}(\cdot)$ performs classification. For each input traffic sample $x_i$, its penultimate-layer representation is $z_i=g_{\theta^{(3)}}(x_i)\in\mathbb{R}^{M}$, where $M$ denotes the feature dimension.

In Phase 2, unknown traffic detection is performed in the feature space rather than relying on softmax probabilities. This avoids the overconfidence problem of neural networks on unknown samples and enables the model to characterize the geometric support regions of known encrypted traffic.

For the $l$-th known class, we collect the features of training samples with observed label $l$ as
\begin{equation}
	\mathcal{Z}_{l}
	=
	\left\{
	z_i=g_{\theta^{(3)}}(x_i)
	\mid
	\tilde{y}_i=l
	\right\}.
\end{equation}
These features are arranged column-wise into $X_l=[z_{l,1},z_{l,2},\ldots,z_{l,N_l}]\in\mathbb{R}^{M\times N_l}$, where $N_l$ is the number of samples in class $l$. We then apply Singular Value Decomposition:
\begin{equation}
	X_l=U_l\Sigma_l V_l^{T},
\end{equation}
where $U_l$ contains the left singular vectors and $r_l=\mathrm{rank}(X_l)$.

Given a global subspace dimension parameter $k$, we retain $k_l=\min(k,r_l)$ principal directions for class $l$ and construct its subspace basis as $U_{l,k_l}=[u_{l,1},u_{l,2},\ldots,u_{l,k_l}]\in\mathbb{R}^{M\times k_l}$. The corresponding class subspace is defined as $\mathcal{S}_{l,k_l}=\mathrm{span}(U_{l,k_l})$. Therefore, the known-class support region in the feature space is represented by the union of all class-wise subspaces:
\begin{equation}
	\mathcal{S}_{in}
	=
	\bigcup_{l=1}^{c}\mathcal{S}_{l,k_l}.
\end{equation}
Compared with a single-direction representation, the multidimensional subspace preserves multiple dominant variation directions within each class, making it more suitable for modeling encrypted traffic variations caused by different sessions, behaviors, and statistical patterns.

Given a test sample $x$, its deep feature is $z=g_{\theta^{(3)}}(x)\in\mathbb{R}^{M}$. To measure its geometric consistency with the $l$-th known-class subspace, we compute the normalized projection similarity:
\begin{equation}
	\rho_l(z)
	=
	\frac{\|U_{l,k_l}^{T}z\|_2}{\|z\|_2},
	\quad
	l=1,2,\ldots,c.
	\label{eq:subspace_similarity}
\end{equation}
Since $U_{l,k_l}$ is composed of orthonormal singular vectors, $0\leq\rho_l(z)\leq1$. A larger $\rho_l(z)$ indicates that the test sample is closer to the $l$-th known-class subspace.

The maximum similarity to all known-class subspaces is defined as
\begin{equation}
	\rho_{\max}(z)
	=
	\max_{1\leq l\leq c}
	\frac{\|U_{l,k_l}^{T}z\|_2}{\|z\|_2}.
	\label{eq:max_similarity}
\end{equation}
Based on this, the multi-subspace angular deviation is defined as
\begin{equation}
	\delta_k(z)
	=
	\arccos(\rho_{\max}(z))
	=
	\arccos
	\left(
	\max_{1\leq l\leq c}
	\frac{\|U_{l,k_l}^{T}z\|_2}{\|z\|_2}
	\right),
	\label{eq:angular_deviation}
\end{equation}
where $\delta_k(z)\in[0,\pi/2]$. A smaller $\delta_k(z)$ means that the sample is close to one known-class subspace and is more likely to be known traffic. Conversely, a larger $\delta_k(z)$ indicates that the sample deviates from all known-class subspaces and is more likely to be unknown traffic. Thus, $\delta_k(z)$ is used as the unknown detection score.

\subsubsection{Test-Time Unknown Traffic Detection}

For a test sample accepted as known traffic, its class is determined by the nearest known-class subspace:
\begin{equation}
	\hat{y}_{in}(x)
	=
	\arg\max_{1\leq l\leq c}\rho_l(z).
	\label{eq:known_class_prediction}
\end{equation}
This prediction strategy relies on the geometric relationship between the sample feature and class subspaces, rather than the softmax output, thereby reducing the impact of softmax overconfidence in open environments.

When $k=1$, each class subspace is reduced to the first principal singular direction, i.e., $U_{l,1}=u_{l,1}$, and $\|U_{l,1}^{T}z\|_2=|u_{l,1}^{T}z|$. In this case, the angular deviation becomes
\begin{equation}
	\delta_1(z)
	=
	\arccos
	\left(
	\max_{1\leq l\leq c}
	\frac{|u_{l,1}^{T}z|}{\|z\|_2}
	\right).
	\label{eq:single_direction_deviation}
\end{equation}
This corresponds to a single-direction spectral discrepancy measure. In contrast, our method uses multiple principal directions for each class, i.e., $\mathcal{S}_{l,1}\subseteq\mathcal{S}_{l,k_l}$ when $k_l>1$, which better captures complex intra-class variations of encrypted traffic and reduces the false rejection of known samples.

The parameter $k$ controls the coverage of known-class subspaces. A smaller $k$ produces more compact subspaces and stricter unknown rejection, while a larger $k$ covers more intra-class variations but may increase the chance of accepting unknown samples. Therefore, $k$ balances known-class coverage and unknown-class rejection.

Since unknown traffic is unavailable during training in practical open environments, the decision threshold is estimated using only known-class training samples. Specifically, after constructing the class-wise subspaces, we compute the angular deviation scores $\{\delta_k(z_i)\}_{i=1}^{N}$ for the training samples and set the threshold as
\begin{equation}
	\tau
	=
	Q_q\left(\{\delta_k(z_i)\}_{i=1}^{N}\right),
	\label{eq:threshold}
\end{equation}
where $Q_q(\cdot)$ denotes the $q$-quantile. A larger $q$ yields a looser threshold and accepts more samples as known, while a smaller $q$ gives a stricter threshold and rejects more deviated samples as unknown.

For a test sample $x$, the unknown traffic decision rule is
\begin{equation}
	\hat{o}(x)
	=
	\begin{cases}
		0, & \delta_k(z)\leq \tau,\\
		1, & \delta_k(z)>\tau,
	\end{cases}
	\label{eq:unknown_decision}
\end{equation}
where $\hat{o}(x)=0$ indicates known-class traffic and $\hat{o}(x)=1$ indicates unknown-class traffic.




%\begin{figure*}[ht]
%	\centering
%	\includegraphics[width=0.45\textwidth]{img/visual}
%	\setlength{\abovecaptionskip}{0cm}
%	\caption{Different visual.}
%	\label{visual}
%	\vspace{-0.2cm}
%\end{figure*}






\section{Experimental Evaluation and Analysis}\label{Experiment}


%We first introduce the datasets, noise settings, and training details used in our experiments. Then, under different noise types and noise ratios, we compare the proposed method with existing label-noise handling methods and unknown traffic detection methods. Finally, ablation studies and parameter sensitivity analyses are conducted to further verify the effectiveness of each component and the rationality of key parameter settings. Specifically, our experiments are designed to answer the following three research questions:

To validate the effectiveness of the proposed model, we conduct a series of experiments to comprehensively evaluate its known-class encrypted traffic classification capability under noisy-label conditions and its unknown attack detection capability in open-world environments. Specifically, our experiments are designed to answer the following four research questions:


\begin{itemize}
	\item \textbf{RQ1: How does our model perform in known-class encrypted traffic classification under noisy labels?}
	\item \textbf{RQ2: How does our model perform in an open-world environment with unknown attacks?}
	\item \textbf{RQ3: How effective are the proposed components and parameter settings?}
	\item \textbf{RQ4: How well does our model generalize to different datasets under noisy-label conditions?}
\end{itemize}

\subsection{Datasets}

We use two representative and relatively recent network attack detection datasets, namely MAL\_TLS2023 and TII-SSRC-23. The MAL\_TLS2023 dataset contains traffic generated by 23 malware families, including Qakbot, PandaZeuS, Tiggre, and others. All traffic flows are encrypted using TLS 1.2\footnote{\url{https://github.com/gcx-Yuan/BoAu}}. After preprocessing, the dataset contains 92,034 network flows.

The TII-SSRC-23 dataset contains 32 classes of benign and attack traffic\footnote{\url{https://www.kaggle.com/datasets/daniaherzalla/tii-ssrc-23}}. Among them, 6 classes correspond to benign traffic, such as Audio, Background, Text, and Video, while the remaining 26 classes correspond to malicious traffic. These malicious classes cover four major attack categories: Bruteforce, DoS, Botnet, and Information Gathering. After preprocessing, this dataset contains 8,656,767 network flows, making it one of the public datasets with the most diverse attack types.

\subsection{Experimental Settings}

\subsubsection{Noise Settings}


To simulate label noise in real-world traffic annotation, we consider two noisy-label settings: symmetric and asymmetric noise \cite{CoLD}. Let $\eta$ denote the noise ratio, i.e., the proportion of training samples selected for label perturbation.

In the symmetric setting, each selected label is uniformly replaced by one of the other $c-1$ classes. This setting assumes that annotation errors are random and independent of class semantics or traffic-pattern similarities.

Asymmetric noise better reflects realistic mislabeling in network attack detection, where classes with similar protocols, attack behaviors, or statistical patterns are more likely to be confused. Therefore, labels are flipped only between semantically or behaviorally similar classes. For example, in TII-SSRC-23, different Mirai DDoS subtypes share similar attack sources and objectives but differ in protocol or packet construction. Thus, Mirai DDoS UDP or Mirai DDoS HTTP may be mislabeled as Mirai DDoS DNS.

%To simulate the label noise commonly observed in real-world network traffic annotation, we design two types of noisy-label settings: symmetric label noise and asymmetric label noise. Let the noise ratio be denoted as $\eta$, which indicates that a proportion $\eta$ of samples are randomly selected from the training set for label perturbation.
%
%In the symmetric noise setting, the labels of the selected samples are randomly replaced by any other class except their original class. Specifically, for a sample with true label $y$, its corrupted label is uniformly sampled from the remaining $c-1$ classes. This setting assumes that annotation errors are completely random and independent of semantic relationships or traffic-pattern similarities among classes.
%
%In contrast, asymmetric noise is more consistent with the annotation errors encountered in realistic network attack detection scenarios. Since some network traffic classes are similar in terms of protocols, attack behaviors, traffic characteristics, or statistical distributions, human annotators or automatic labeling systems are more likely to confuse such closely related classes. Therefore, in the asymmetric noise setting, we do not randomly replace labels with arbitrary classes. Instead, label flipping or one-way replacement is performed only among classes with similar semantics or similar traffic patterns. For example, in the TII-SSRC-23 dataset, multiple DDoS subtypes of Mirai botnet attacks share similar attack sources and objectives, differing mainly in specific protocols or packet construction strategies. Thus, Mirai DDoS UDP may be mislabeled as Mirai DDoS DNS, and Mirai DDoS HTTP may also be mislabeled as Mirai DDoS DNS.


\subsubsection{Experimental Environment}

To evaluate the capability of the Argus model in detecting unknown malicious traffic, we construct known-class and unknown-class splits on both datasets. In addition, unknown traffic from external datasets is introduced to enhance the realism and difficulty of the test scenarios.

For MAL\_TLS2023, 21 of the 23 malware traffic classes are randomly selected as known classes for training, while the remaining 2 classes are treated as internal unknown classes and used only for testing. In addition, 5,000 flows from CipherSpectrum~\cite{CipherSpectrum} are used as external unknown traffic. These TLS 1.3 encrypted web flows are processed with the same pipeline as MAL\_TLS2023. Thus, the MAL\_TLS2023 test set contains 10,000 unknown flows in total, including both internal unknown traffic and external CipherSpectrum traffic.

For TII-SSRC-23, 20 of the 32 subclasses, covering both benign and malicious traffic, are used as known classes for training and closed-set evaluation. The remaining 12 DoS-related classes are used as internal unknown classes to assess detection of unseen attacks. Moreover, 5,000 flows from 9 attack categories in CIC-UNSW-NB15 \cite{CIC-UNSW-NB15}, including Fuzzers, Analysis, Reconnaissance, Shellcode, and Worms, are introduced as external unknown traffic. Consequently, the TII-SSRC-23 test set also contains 10,000 unknown flows, comprising internal unknown DoS traffic and external CIC-UNSW-NB15 attack traffic.

\subsubsection{Training Details}

To ensure fair comparison among different methods, we use the same data splits, noise types, and noise ratios in all experiments. The proposed method adopts a CNN-based deep residual neural network as the backbone model to extract discriminative features from network traffic. During training, the confidence threshold is set to $\gamma=0.1$, the batch size is set to 256, and the dimension of the deep feature representation is set to $M=256$. The overall training process consists of three stages. Stage 1 is trained for 240 epochs, where the first 130 epochs use the complementary-label loss. Stage 2 is trained for 240 epochs. Stage 3 is trained for 200 epochs. All experimental results are averaged over at least five independent runs.

\subsection{Comparison Methods}
We compare the proposed method with various state-of-the-art methods. The compared methods mainly fall into two categories: (i) label-noise handling methods, including AEGIS-Net, RAPIER, MCRe, DSDIR, and MORSE; and (ii) unknown traffic detection methods, including FOSS, RFG-HELAD, OpenMax, as well as several commonly used unknown detection heads, i.e., MSP, Energy, and Cosine.

%\begin{table}[t]
%	\footnotesize
%	\centering
%	\caption{Summary of the selected comparison methods.}
%	\label{tab:comparison_methods}
%	\begin{tabular}{ll}
%		\toprule
%		\textbf{Category} & \textbf{Methods} \\
%		\midrule
%		Label-noise handling methods 
%		& AEGIS-Net, RAPIER, MCRe, DSDIR, MORSE \\
%		Unknown traffic detection methods 
%		& FOSS, RFG-HELAD, OpenMax, MSP, Energy, Cosine \\
%		\bottomrule
%	\end{tabular}
%	\vspace{-0.4cm}
%\end{table}

\begin{itemize}
	\item AEGIS-Net \cite{AEGIS-Net} detects and corrects mislabeled samples by measuring their similarity to multiple class prototypes.
	
	\item RAPIER \cite{RAPIER} selects clean samples using Masked Autoencoder for Distribution Estimation (MADE) based on sample density characteristics.
	
	\item MCRe \cite{MCRe} combines data cleaning and robust training by imposing information integrity, cluster separability, and core proximity constraints.
	
	\item DSDIR \cite{DSDIR} addresses noisy labels and long-tailed distributions through noise-independent clean sample selection.
	
	\item MORSE \cite{MORSE} mitigates high label noise and class imbalance by combining semi-supervised learning with sample reweighting.
	
	\item FOSS \cite{FOSS} detects unknown traffic using an isolation-forest-based multi-tree model and deviation-based anomaly scoring.
	
	\item RFG-HELAD \cite{RFG-HELAD} employs heterogeneous ensemble learning with $K$-class and $(K+1)$-class submodels to distinguish known and unknown attacks.
	
	\item OpenMax \cite{OpenMax} applies Extreme Value Theory (EVT) to model distances from known-class centers for open-set recognition.
\end{itemize}

\subsection{Known Traffic Classification Experiments (RQ1)}


\begin{table*}[ht]
	\footnotesize
	\begin{center}
		\caption{The closed-world performance of all methods under different noise ratios on the MAL\_TLS2023 dataset.}
		\label{closed_world_MAL_TLS2023}
		\setlength{\tabcolsep}{5pt}
		\renewcommand{\arraystretch}{1.2}
		\begin{threeparttable}
			\resizebox{\textwidth}{!}{
				\begin{tabular}{lccccccccccccccccc}
					\toprule
					\textbf{Noise Type} 
					& $|$ 
					& \multicolumn{2}{c}{\textbf{None}} 
					& $|$ 
					& \multicolumn{6}{c}{\textbf{Symmetric}} 
					& $|$ 
					& \multicolumn{6}{c}{\textbf{Asymmetric}} \\
					\midrule
					
					\textbf{Noise Ratio} 
					& $|$ 
					& \multicolumn{2}{c}{\textbf{0\%}} 
					& $|$ 
					& \multicolumn{2}{c}{\textbf{10\%}} 
					& \multicolumn{2}{c}{\textbf{30\%}} 
					& \multicolumn{2}{c}{\textbf{50\%}} 
					& $|$ 
					& \multicolumn{2}{c}{\textbf{10\%}} 
					& \multicolumn{2}{c}{\textbf{30\%}} 
					& \multicolumn{2}{c}{\textbf{50\%}} \\
					\midrule
					
					\textbf{Methods} 
					& $|$ 
					& \textbf{Acc} & \textbf{F1} 
					& $|$ 
					& \textbf{Acc} & \textbf{F1} 
					& \textbf{Acc} & \textbf{F1} 
					& \textbf{Acc} & \textbf{F1} 
					& $|$ 
					& \textbf{Acc} & \textbf{F1} 
					& \textbf{Acc} & \textbf{F1} 
					& \textbf{Acc} & \textbf{F1} \\
					\midrule
					
					AEGIS-Net 
					& $|$ 
					& 97.23 & 97.61 
					& $|$ 
					& 95.92 & 95.85 
					& 95.53 & 95.61 
					& 94.83 & 94.85 
					& $|$ 
					& 92.34 & 92.57 
					& \textbf{91.72} & \textbf{91.47} 
					& 87.74 & 87.36 \\
					
					RAPIER 
					& $|$ 
					& 91.62 & 91.43 
					& $|$ 
					& 88.40 & 88.32 
					& 86.39 & 86.58 
					& 83.65 & 83.49 
					& $|$ 
					& 91.72 & 91.56 
					& 85.69 & 85.77 
					& 82.84 & 82.91 \\
					
					MCRe 
					& $|$ 
					& 96.29 & 96.32 
					& $|$ 
					& 94.54 & 94.43 
					& 93.57 & 93.35 
					& 91.38 & 91.30 
					& $|$ 
					& 92.56 & 92.48 
					& 90.77 & 90.68 
					& \textbf{89.23} & 89.19 \\
					
					DSDIR 
					& $|$ 
					& 95.75 & 95.88 
					& $|$ 
					& 94.48 & 94.51 
					& 93.61 & 93.63 
					& 92.14 & 92.35 
					& $|$ 
					& 90.75 & 90.43 
					& 88.58 & 88.60 
					& 86.23 & 86.26 \\
					
					MORSE 
					& $|$ 
					& 89.58 & 89.43 
					& $|$ 
					& 87.35 & 87.40 
					& 86.62 & 86.59 
					& 84.81 & 84.70 
					& $|$ 
					& 88.58 & 88.49 
					& 86.42 & 86.34 
					& 83.45 & 83.72 \\
					
					RFG-HELAD 
					& $|$ 
					& 87.23 & 84.60 
					& $|$ 
					& 86.10 & 82.77 
					& 87.95 & 84.57 
					& 86.90 & 83.54 
					& $|$ 
					& 84.84 & 81.07 
					& 81.77 & 77.16 
					& 59.90 & 57.17 \\
					\midrule
					
					\textbf{Argus} 
					& $|$ 
					& \textbf{98.02} & \textbf{98.01} 
					& $|$ 
					& \textbf{97.47} & \textbf{97.42} 
					& \textbf{97.30} & \textbf{97.24} 
					& \textbf{96.17} & \textbf{96.14} 
					& $|$ 
					& \textbf{93.22} & \textbf{93.81} 
					& 86.04 & 86.71 
					& 88.78 & \textbf{89.55} \\
					
					\bottomrule
				\end{tabular}
			}
		\end{threeparttable}
	\end{center}
	\vspace{-0.7cm}
\end{table*}


\begin{table*}[ht]
	\footnotesize
	\begin{center}
		\caption{The closed-world performance of all methods under different noise ratios on the TII-SSRC-23 dataset.}
		\label{closed_world_TII_SSRC_23}
		\setlength{\tabcolsep}{5pt}
		\renewcommand{\arraystretch}{1.2}
		\begin{threeparttable}
			\resizebox{\textwidth}{!}{
				\begin{tabular}{lccccccccccccccccc}
					\toprule
					\textbf{Noise Type}
					& $|$
					& \multicolumn{2}{c}{\textbf{None}}
					& $|$
					& \multicolumn{6}{c}{\textbf{Symmetric}}
					& $|$
					& \multicolumn{6}{c}{\textbf{Asymmetric}} \\
					\midrule
					
					\textbf{Noise Ratio}
					& $|$
					& \multicolumn{2}{c}{\textbf{0\%}}
					& $|$
					& \multicolumn{2}{c}{\textbf{10\%}}
					& \multicolumn{2}{c}{\textbf{30\%}}
					& \multicolumn{2}{c}{\textbf{50\%}}
					& $|$
					& \multicolumn{2}{c}{\textbf{10\%}}
					& \multicolumn{2}{c}{\textbf{30\%}}
					& \multicolumn{2}{c}{\textbf{50\%}} \\
					\midrule
					
					\textbf{Methods}
					& $|$
					& \textbf{Acc} & \textbf{F1}
					& $|$
					& \textbf{Acc} & \textbf{F1}
					& \textbf{Acc} & \textbf{F1}
					& \textbf{Acc} & \textbf{F1}
					& $|$
					& \textbf{Acc} & \textbf{F1}
					& \textbf{Acc} & \textbf{F1}
					& \textbf{Acc} & \textbf{F1} \\
					\midrule
					
					AEGIS-Net
					& $|$
					& 92.44 & 92.63
					& $|$
					& 92.40 & 92.35
					& 89.68 & 89.65
					& 85.47 & 85.54
					& $|$
					& 90.31 & 90.24
					& 88.69 & 88.65
					& 85.01 & 84.80 \\
					
					RAPIER
					& $|$
					& 90.10 & 90.01
					& $|$
					& 89.95 & 89.77
					& 86.25 & 86.46
					& 82.68 & 82.70
					& $|$
					& 88.91 & 88.85
					& 84.72 & 84.70
					& 81.25 & 81.39 \\
					
					MCRe
					& $|$
					& 89.56 & 89.12
					& $|$
					& 87.29 & 87.31
					& 86.52 & 86.63
					& 84.56 & 84.53
					& $|$
					& 86.55 & 86.38
					& 83.47 & 83.12
					& 79.59 & 79.35 \\
					
					DSDIR
					& $|$
					& 91.23 & 91.19
					& $|$
					& 89.53 & 89.29
					& 87.69 & 87.41
					& 83.86 & 83.09
					& $|$
					& 85.76 & 85.82
					& 82.27 & 82.11
					& 75.29 & 75.41 \\
					
					MORSE
					& $|$
					& 87.33 & 87.30
					& $|$
					& 85.45 & 85.45
					& 82.96 & 82.82
					& 78.23 & 78.12
					& $|$
					& 83.34 & 83.23
					& 81.11 & 81.02
					& 77.00 & 77.04 \\
					
					RFG-HELAD
					& $|$
					& 92.65 & 92.41
					& $|$
					& 92.69 & 92.45
					& 92.62 & 92.43
					& 92.76 & 92.56
					& $|$
					& \textbf{92.70} & \textbf{92.44}
					& \textbf{92.56} & 92.41
					& 85.30 & 83.19 \\
					\midrule
					
					\textbf{Argus}
					& $|$
					& \textbf{93.84} & \textbf{93.89}
					& $|$
					& \textbf{94.63} & \textbf{94.75}
					& \textbf{94.37} & \textbf{94.35}
					& \textbf{94.84} & \textbf{94.88}
					& $|$
					& 91.02 & 92.42
					& 92.16 & \textbf{93.18}
					& \textbf{91.45} & \textbf{92.74} \\
					
					\bottomrule
				\end{tabular}
			}
		\end{threeparttable}
	\end{center}
	\vspace{-0.4cm}
\end{table*}


In this subsection, we evaluate the known-class traffic classification performance of the proposed Argus model under noisy-label scenarios on MAL\_TLS2023 and TII-SSRC-23. We consider both symmetric and asymmetric noise with noise ratios from 0\% to 50\%, and compare our method with six representative baselines. The results are reported in Tables~\ref{closed_world_MAL_TLS2023} and~\ref{closed_world_TII_SSRC_23}.

\textbf{1) Impact of noise type and noise ratio.}
As the noise ratio increases, most methods suffer from performance degradation, but their robustness differs significantly. On MAL\_TLS2023, Argus achieves the best clean-setting performance, with 98.02\% accuracy and 98.01\% F1-score. Under symmetric noise, it still maintains 96.17\% accuracy and 96.14\% F1-score at a 50\% noise ratio, clearly outperforming RAPIER, MORSE, and RFG-HELAD. This demonstrates its strong robustness to severe random label corruption. Under asymmetric noise, where mislabeled samples usually occur between semantically similar or boundary-ambiguous classes, classification becomes more challenging. Nevertheless, Argus still achieves the best result under 10\% asymmetric noise on MAL\_TLS2023, and similar trends are observed on TII-SSRC-23.

\textbf{2) Cross-dataset performance.}
Across the two datasets, Argus shows stable and competitive performance. On MAL\_TLS2023, AEGIS-Net performs well under some asymmetric noise settings, especially at 30\% asymmetric noise, possibly because its multi-prototype mechanism captures local intra-class structures. However, our method performs better in symmetric and low-ratio asymmetric noise settings, indicating that complementary label learning and confidence-based sample selection more effectively suppress noisy-label interference. On TII-SSRC-23, the advantage of our method is more evident. In contrast, RFG-HELAD exhibits noticeable fluctuations and degrades under 50\% asymmetric noise, likely because it relies on relatively clean labels to learn reliable decision boundaries. By dynamically separating high-confidence and low-confidence samples, our method reduces the accumulated negative impact of noisy labels and achieves more stable performance.

\textbf{3) Robustness to label noise.}
Overall, Argus achieves the best or highly competitive accuracy and F1-score in most settings, especially under symmetric and high-ratio noise. Its robustness mainly comes from two designs. First, complementary label learning avoids directly fitting potentially incorrect labels at the early training stage and guides the model by excluding unlikely classes. Second, confidence-based fitting exploits reliable samples while reducing the influence of low-confidence noisy ones. These two mechanisms jointly form a ``noise suppression--confidence enhancement'' training process, leading to robust known-class encrypted traffic classification.

Figure~\ref{fig:acc} further shows that the samples used for final retraining in Stage~3 maintain high label accuracy. This verifies that the first two stages can effectively filter noisy samples and improve training-data purity. The label accuracy under symmetric noise is generally higher than that under asymmetric noise, since randomly corrupted labels are easier to distinguish by confidence differences, whereas asymmetric noise often occurs between similar traffic classes and is therefore harder to remove. As the noise ratio increases, the label accuracy of Stage~3 data decreases, indicating that high-noise settings remain challenging. Nevertheless, the results confirm that Argus consistently improves training-data quality under different noisy-label conditions.

\begin{figure*}[ht]
	\centering
	\begin{minipage}[t]{0.5\textwidth}
		\centering
		\includegraphics[width=\linewidth]{img/acc}
		\captionof{figure}{The label accuracy of the data used for model training in Stage 3 under different noise settings on the MAL\_TLS2023 and TII-SSRC-23 datasets.}
		\label{fig:acc}
	\end{minipage}
	\hspace{0.0\textwidth}
	\begin{minipage}[t]{0.46\textwidth}
		\centering
		\includegraphics[width=\linewidth]{img/AUC}
		\captionof{figure}{ROC curves of Argus on the MAL\_TLS2023 and TII-SSRC-23 datasets.}
		\label{fig:AUC}
	\end{minipage}
\end{figure*}


\begin{table*}[ht] 
	\footnotesize
	\begin{center} 
		\caption{The open-world performance of all methods under different noise ratios on the MAL\_TLS2023 and TII-SSRC-23 datasets.} 
		\label{open_world_noise_performance} 
		\setlength{\tabcolsep}{4pt} 
		\renewcommand{\arraystretch}{1.15} 
		\begin{threeparttable} 
			\resizebox{\textwidth}{!}{ 
				\begin{tabular}{clccccccccccccccccc} 
					\toprule
					
					\multirow{3}{*}{\textbf{Dataset}} 
					& \textbf{Noise Type} 
					& $|$
					& \multicolumn{2}{c}{\textbf{None}} 
					& $|$
					& \multicolumn{6}{c}{\textbf{Symmetric}} 
					& $|$
					& \multicolumn{6}{c}{\textbf{Asymmetric}} \\ 
					\cmidrule(lr){2-19} 
					
					& \textbf{Noise Ratio} 
					& $|$
					& \multicolumn{2}{c}{\textbf{0\%}} 
					& $|$
					& \multicolumn{2}{c}{\textbf{10\%}} 
					& \multicolumn{2}{c}{\textbf{30\%}} 
					& \multicolumn{2}{c}{\textbf{50\%}} 
					& $|$
					& \multicolumn{2}{c}{\textbf{10\%}} 
					& \multicolumn{2}{c}{\textbf{30\%}} 
					& \multicolumn{2}{c}{\textbf{50\%}} \\ 
					\cmidrule(lr){2-19} 
					
					& \textbf{Methods} 
					& $|$
					& \textbf{Acc} & \textbf{F1} 
					& $|$
					& \textbf{Acc} & \textbf{F1} 
					& \textbf{Acc} & \textbf{F1} 
					& \textbf{Acc} & \textbf{F1} 
					& $|$
					& \textbf{Acc} & \textbf{F1} 
					& \textbf{Acc} & \textbf{F1} 
					& \textbf{Acc} & \textbf{F1} \\ 
					\midrule
					
					\multirow{8}{*}{\rotatebox{90}{MAL\_TLS2023}} 
					& Argus-MSP
					& $|$
					& 71.31 & 80.55
					& $|$
					& 69.64 & 79.67
					& 59.48 & 74.58
					& 70.69 & 80.23
					& $|$
					& 65.72 & 77.61
					& 62.03 & 75.82
					& 61.55 & 75.51 \\ 
					
					& Argus-Energy
					& $|$
					& 78.16 & 84.50
					& $|$
					& 78.04 & 84.40
					& 76.93 & 83.72
					& 70.93 & 80.36
					& $|$
					& 77.03 & 83.79
					& 77.11 & 83.84
					& 77.04 & 83.78 \\ 
					
					& Argus-Cosine
					& $|$
					& 93.82 & 94.40
					& $|$
					& 93.82 & 94.40
					& 93.73 & 94.32
					& 92.12 & \textbf{93.06}
					& $|$
					& 93.84 & 94.41
					& 93.83 & \textbf{94.40}
					& 93.83 & 94.41 \\ 
					
					& AEGIS-Net
					& $|$
					& 93.81 & 94.39
					& $|$
					& 93.78 & 94.36
					& 93.67 & 94.27
					& 92.06 & 93.01
					& $|$
					& 87.08 & 89.29
					& 93.74 & 94.33
					& 93.74 & 94.33 \\ 
					
					& Openmax
					& $|$
					& 93.87 & 92.98
					& $|$
					& 92.01 & 90.46
					& 90.98 & 88.98
					& 91.43 & 89.66
					& $|$
					& 88.46 & 83.41
					& 84.54 & 76.47
					& 78.21 & 63.45 \\ 
					
					& FOSS
					& $|$
					& 64.63 & 77.10
					& $|$
					& 60.23 & 74.96
					& 61.11 & 75.38
					& 62.57 & 76.08
					& $|$
					& 62.95 & 76.27
					& 62.17 & 75.88
					& 62.17 & 75.88 \\ 
					
					& RFG-HELAD
					& $|$
					& 94.23 & 94.23
					& $|$
					& 91.19 & 91.21
					& 91.18 & 91.20
					& 91.15 & 91.16
					& $|$
					& 91.10 & 91.11
					& 91.18 & 91.20
					& 91.09 & 91.11 \\ 
					
					& \textbf{Ours} 
					& $|$
					& \textbf{96.49} & \textbf{95.51} 
					& $|$
					& \textbf{96.01} & \textbf{94.87} 
					& \textbf{96.56} & \textbf{95.59} 
					& \textbf{94.02} & 92.07 
					& $|$
					& \textbf{96.56} & \textbf{95.60} 
					& \textbf{95.55} & 94.24 
					& \textbf{95.52} & \textbf{94.20} \\ 
					\midrule
					
					\multirow{8}{*}{\rotatebox{90}{TII-SSRC-23}} 
					& Argus-MSP
					& $|$
					& 85.18 & 88.53
					& $|$
					& 88.99 & 91.21
					& 85.54 & 88.79
					& 71.45 & 80.05
					& $|$
					& 63.42 & 75.80
					& 63.80 & 76.00
					& 62.29 & 75.23 \\ 
					
					& Argus-Energy
					& $|$
					& 86.97 & 89.76
					& $|$
					& 88.75 & 91.03
					& 83.18 & 87.20
					& 82.44 & 86.70
					& $|$
					& 85.89 & 89.02
					& 65.43 & 76.80
					& 60.66 & 74.41 \\ 
					
					& Argus-Cosine
					& $|$
					& 90.82 & 91.47
					& $|$
					& 90.63 & \textbf{91.32}
					& 91.16 & \textbf{91.75}
					& 91.16 & \textbf{91.75}
					& $|$
					& 90.67 & 90.35
					& 91.14 & 91.74
					& 91.23 & 91.81 \\ 
					
					& AEGIS-Net
					& $|$
					& 91.71 & 93.24
					& $|$
					& 89.35 & 90.98
					& 89.43 & 91.04
					& 89.81 & 91.35
					& $|$
					& 90.77 & 92.05
					& 89.91 & 91.40
					& 90.85 & 92.11 \\ 
					
					& Openmax
					& $|$
					& 83.32 & 75.80
					& $|$
					& 77.97 & 65.32
					& 82.16 & 73.66
					& 75.57 & 59.92
					& $|$
					& 83.20 & 75.56
					& 85.81 & 80.08
					& 64.53 & 28.94 \\ 
					
					& FOSS
					& $|$
					& 75.49 & 82.36
					& $|$
					& 72.34 & 80.55
					& 66.77 & 77.51
					& 67.86 & 78.08
					& $|$
					& 74.00 & 81.50
					& 68.98 & 78.69
					& 71.90 & 80.30 \\ 
					
					& RFG-HELAD
					& $|$
					& 79.70 & 78.55
					& $|$
					& 86.70 & 86.41
					& 82.49 & 81.77
					& 88.20 & 88.00
					& $|$
					& 78.75 & 77.42
					& 79.40 & 78.20
					& 80.20 & 79.14 \\ 
					
					& \textbf{Argus} 
					& $|$
					& \textbf{94.76} & \textbf{93.89} 
					& $|$
					& \textbf{92.54} & 90.50 
					& \textbf{91.82} & 89.50 
					& \textbf{91.94} & 89.63 
					& $|$
					& \textbf{94.01} & \textbf{92.52} 
					& \textbf{95.34} & \textbf{94.28} 
					& \textbf{93.80} & \textbf{92.24} \\ 
					
					\bottomrule
				\end{tabular} 
			}
		\vspace{2pt}
		\begin{tablenotes}[flushleft]
			\footnotesize
			\item[] 
			\parbox{\textwidth}{
				\textit{Note:} 
				Argus-MSP uses the Maximum Softmax Probability of the classifier's logits to measure detection confidence~\cite{MSP}.
				Argus-Energy calculates the Energy score based on the LogSumExp of the output logits~\cite{Energy}.
				Argus-Cosine measures the maximum cosine similarity between the normalized test sample feature and the training feature bank~\cite{Cosine}.
			}
		\end{tablenotes}
		\end{threeparttable} 
	\end{center} 
	\vspace{-0.6cm} 
\end{table*}

% Note

\subsection{Unknown Traffic Detection Experiments (RQ2)}

To evaluate the effectiveness of Argus in detecting unknown malicious traffic, we conduct experiments under different noise settings. Based on the results in Table~\ref{open_world_noise_performance}, we make the following observations.

\textbf{1) The ``denoising first and detection later'' design effectively supports unknown traffic recognition.}
Argus achieves the best detection Acc and competitive F1-score in most settings on both MAL\_TLS2023 and TII-SSRC-23. On MAL\_TLS2023, Argus consistently outperforms OpenMax, FOSS, and RFG-HELAD in terms of Acc under clean labels, symmetric noise, and asymmetric noise. On TII-SSRC-23, Argus maintains Acc above 91\% across all noise settings and reaches 95.34\% under 30\% asymmetric noise. These results show that the three-stage denoising training in Phase 1 can learn cleaner and more stable traffic representations from noisy data, providing a reliable feature basis for unknown traffic detection in Phase 2. By suppressing early fitting to incorrect labels through complementary label learning and further refining the model with high-confidence samples, Argus produces a clearer separation between unknown samples and the feature support regions of known classes.

\textbf{2) The multidimensional subspace angular deviation provides more stable detection.}
Compared with Ours-MSP and Ours-Energy, which show unstable performance in several noisy scenarios, Argus achieves more robust results. This indicates that detection scores based on maximum softmax probability or energy are easily affected by model overconfidence and noisy-label contamination. Ours-Cosine performs better, suggesting that feature-space similarity is more reliable than classifier confidence. However, cosine similarity still relies on a single-direction or point-wise comparison, which is insufficient to describe the complex intra-class variations of encrypted traffic. In contrast, Argus constructs a multidimensional principal subspace for each known class and measures unknownness by the angular deviation between a test sample and all known-class subspaces. This design better captures variations caused by different session states and communication behaviors, reducing the false rejection of known-class variants while effectively detecting samples that deviate from all known subspaces.

\textbf{3) Argus remains robust under high-ratio and asymmetric label noise.}
Asymmetric noise is more challenging than random symmetric noise because it usually occurs between semantically similar classes or classes with similar traffic patterns, which can seriously distort known-class boundaries. As shown in Table~\ref{open_world_noise_performance}, traditional methods degrade noticeably under such conditions. For example, OpenMax drops to only 28.94\% F1 on TII-SSRC-23 under 50\% asymmetric noise, and FOSS and RFG-HELAD also fail to maintain stable performance. This is mainly because these methods do not explicitly handle label noise, making their detection boundaries vulnerable to mislabeled samples. In comparison, Argus still achieves 93.80\% Acc and 92.24\% F1 under the same setting. The reason is twofold: Stage 3 uses only high-confidence clean samples for robust retraining, reducing the contamination of noisy samples in feature learning and subspace construction; meanwhile, the Phase 2 detector makes decisions based on geometric consistency with known-class feature subspaces rather than softmax outputs. Therefore, Argus can maintain a clear known/unknown boundary even under severe and structured label noise.



\begin{figure*}[ht]
	\centering
	\includegraphics[width=0.9\textwidth]{img/anomaly}
	\setlength{\abovecaptionskip}{0cm}
	\caption{Histogram of the anomaly score $\delta_k$ and the detection threshold in Argus, where samples with scores lower than the threshold are classified as known classes, while those with scores higher than the threshold are classified as unknown classes.}
	\label{anomaly}
	\vspace{-0.2cm}
\end{figure*}


Figure~\ref{fig:AUC} presents the ROC curves and AUC results of Argus on MAL\_TLS2023 and TII-SSRC-23 under different noise settings. As shown in the figure, the ROC curves are consistently close to the upper-left corner and clearly above the random detection baseline, indicating that Argus can effectively distinguish known-class traffic from unknown-class traffic. On MAL\_TLS2023, Argus achieves AUC values above 98.97\% across all noise settings, with the highest AUC of 99.67\% under 10\% asymmetric noise. On TII-SSRC-23, despite the larger number of classes and more complex traffic patterns, Argus still maintains AUC values above 97.32\%. These results demonstrate that the multidimensional subspace angular deviation provides a reliable detection boundary for unknown traffic. Moreover, the consistently high AUC under severe noise conditions shows that the complementary label learning and confidence-based sample selection in Phase~1 effectively reduce noisy-label contamination, enabling the class subspaces constructed in Phase~2 to remain robust.

Figure~\ref{anomaly} further illustrates the anomaly score distributions and detection thresholds of Argus under different asymmetric noise rates. Samples with anomaly scores below the threshold are classified as known classes, while those above the threshold are identified as unknown classes. On MAL\_TLS2023 and TII-SSRC-23, known traffic is mainly concentrated in the low-score region, whereas unknown traffic is distributed in higher-score regions, forming a clear separation around the threshold. This confirms that the proposed anomaly score can effectively measure the geometric consistency between a sample and the known-class feature subspaces: known samples are close to one of the class subspaces and thus obtain lower scores, while unknown samples deviate from all known subspaces and receive higher scores. Even under the high-noise Asym-50\% setting, the two distributions remain well separated, indicating that Phase~1 successfully mitigates the influence of noisy labels on subspace construction. These observations further verify the robustness of Argus and the rationality of performing unknown traffic detection from the perspective of feature-space geometric deviation rather than softmax confidence.

%Figure \ref{AUC}




%\begin{figure*}[ht]
%	\centering
%	\includegraphics[width=0.95\textwidth]{img/acc}
%	\setlength{\abovecaptionskip}{0cm}
%	\caption{Different acc.}
%	\label{acc}
%	\vspace{-0.2cm}
%\end{figure*}

%\begin{figure*}[ht]
%	\centering
%	\includegraphics[width=0.85\textwidth]{img/AUC}
%	\setlength{\abovecaptionskip}{0cm}
%	\caption{Different AUC.}
%	\label{AUC}
%	\vspace{-0.2cm}
%\end{figure*}




\subsection{Ablation Study (RQ3)}
This subsection answers RQ3 by conducting ablation studies and parameter sensitivity analyses to evaluate the effectiveness of the key components in Argus. Specifically, we examine the contribution of each training stage and analyze the impact of two important hyperparameters, $\gamma$ and $k$, on both known-class classification and unknown-traffic detection.


\subsubsection{Effect of each stage}

\begin{figure*}[ht]
	\centering
	\includegraphics[width=0.98\textwidth]{img/ablation}
	\setlength{\abovecaptionskip}{0cm}
	\caption{Ablation results of the three stages from Stage 1 to Stage 3 in Phase 1 of the Argus model.}
	\label{ablation}
	\vspace{-0.4cm}
\end{figure*}



The ablation results in Figure~\ref{ablation} show that each module in the three-stage training strategy of Argus contributes positively to the final performance. Compared with using only Stage 1, incorporating Stage 2 brings slight improvements in most low-noise scenarios, indicating that fitting target labels on high-confidence samples can further exploit reliable supervisory information. However, the most significant performance gain is achieved after introducing Stage 3. By explicitly selecting high-confidence clean samples and continuing model training only on the purified sample subset, Stage 3 effectively reduces the contamination of the final decision boundary and feature representations caused by suspected noisy samples. As a result, the full model achieves the best or near-best performance under both symmetric and asymmetric noise, with more pronounced improvements in high-noise scenarios. In addition, the improvement in unknown traffic detection further demonstrates that the known-class feature representations purified by Stage 3 can construct a more reliable class subspace, thereby enhancing unknown traffic recognition in open-world environments.


\subsubsection{Effect of $\gamma$}
The sensitivity analysis of the confidence threshold $\gamma$ demonstrates that setting $\gamma=0.1$ achieves the best trade-off between known-class classification and unknown-traffic detection. This observation is consistent with the design principle of the proposed multi-stage denoising framework. In Phase 1, complementary label exclusion learning prevents the model from directly fitting potentially corrupted observed labels and instead encourages the formation of robust feature representations through exclusion-based supervision. As training proceeds, clean samples and noisy samples become progressively separable in terms of the prediction confidence assigned to their observed labels. Therefore, a relatively small threshold, i.e., $\gamma=0.1$, is sufficient to retain most reliable samples while filtering out low-confidence noisy instances. Compared with larger thresholds, this setting avoids overly conservative sample selection and preserves a larger amount of informative training data for subsequent positive learning. 

As shown in Figure \ref{gamma}, both the known-class classification accuracy and F1-score reach their highest values at $\gamma=0.1$, achieving 0.9747 and 0.9742, respectively. Meanwhile, the unknown detection module also obtains the best performance, with an accuracy of 0.9664 and an F1-score of 0.9570. When $\gamma$ increases, the performance gradually degrades, especially for $\gamma \geq 0.8$, indicating that an excessively strict confidence threshold discards many useful clean samples and weakens the learned class-discriminative representation. This degradation further affects the construction of known-class feature subspaces in Phase 2, resulting in reduced unknown rejection capability. These results confirm that $\gamma=0.1$ effectively balances noise suppression and data utilization, thereby enabling the proposed framework to learn robust representations for both closed-set classification and open-set unknown detection.

\begin{figure*}[ht]
	\centering
	\begin{minipage}[t]{0.42\textwidth}
		\centering
		\includegraphics[width=\linewidth]{img/gamma}
		\captionof{figure}{The impact of different values of parameter $\gamma$ on model performance under the Sym-10\% setting on the MAL\_TLS2023 dataset.}
		\label{gamma}
	\end{minipage}
	\hspace{0.0\textwidth}
	\begin{minipage}[t]{0.55\textwidth}
		\centering
		\includegraphics[width=\linewidth]{img/k}
		\captionof{figure}{The impact of different values of parameter $k$ on model performance under the Sym-10\% setting on the MAL\_TLS2023 dataset.}
		\label{fig:k}
	\end{minipage}
\end{figure*}

\subsubsection{Effect of $k$}
As shown in Figure~\ref{fig:k}(a) and Figure~\ref{fig:k}(b), the parameter $k$, which controls the maximum dimensionality of each class subspace, has a non-monotonic yet generally saturating effect on model performance. For known-class classification shown in Figure~\ref{fig:k}(a), increasing $k$ from 1 to around 30 consistently improves both Acc and F1. This indicates that a larger subspace dimensionality enables the model to capture richer multimodal intra-class variations in encrypted traffic, thereby improving the robustness of the decision boundary. When $k$ exceeds 30, the performance gradually stabilizes and reaches its peak at $k=50$, suggesting that the learned subspaces have already covered the dominant intra-class structures. Further increasing the dimensionality provides only marginal benefits and may even introduce slight overfitting.

For unknown-class detection shown in Figure~\ref{fig:k}(b), the detection Acc and F1 remain relatively stable when $k$ ranges from 1 to 30, but a local performance drop appears around $k=32$. One possible reason is that an insufficient subspace dimensionality makes the class subspaces overly compact, causing some known-class outliers to be incorrectly rejected as unknown. As $k$ further increases to 50, the model achieves the best overall performance. These results suggest that $k=50$ provides the best trade-off between modeling intra-class diversity and maintaining a clear rejection boundary. It avoids the false rejection of known samples caused by overly narrow subspaces while preventing overly broad subspaces from accepting excessive unknown samples. Therefore, we set $k=50$ in this paper to ensure that the multidimensional principal subspaces can sufficiently model intra-class heterogeneity, thereby improving both closed-world classification robustness and unknown traffic detection capability.

%\begin{figure*}[ht]
%	\centering
%	\includegraphics[width=0.8\textwidth]{img/k}
%	\setlength{\abovecaptionskip}{0cm}
%	\caption{Different $k$.}
%	\label{k}
%	\vspace{-0.3cm}
%\end{figure*}



\subsection{Generalizability Analysis (RQ4)}

\begin{table*}[ht] 
	\footnotesize
	\begin{center} 
		\caption{Classification performance of Argus on different datasets under symmetric noise scenarios.} 
		\label{closed_world_symmetric} 
		\setlength{\tabcolsep}{5pt} 
		\renewcommand{\arraystretch}{1.2} 
		
		\begin{threeparttable} 
			\resizebox{\textwidth}{!}{ 
				\begin{tabular}{lc cc c cc c cc c cc c cc c cc c cc c cc} 
					\toprule
					
					\textbf{Noise Type} 
					& $|$
					& \multicolumn{23}{c}{\textbf{Symmetric}} \\ 
					
					\midrule
					
					\textbf{Noise Ratio} 
					& $|$
					& \multicolumn{2}{c}{\textbf{10\%}} 
					& $|$
					& \multicolumn{2}{c}{\textbf{20\%}} 
					& $|$
					& \multicolumn{2}{c}{\textbf{30\%}} 
					& $|$
					& \multicolumn{2}{c}{\textbf{40\%}} 
					& $|$
					& \multicolumn{2}{c}{\textbf{50\%}} 
					& $|$
					& \multicolumn{2}{c}{\textbf{60\%}} 
					& $|$
					& \multicolumn{2}{c}{\textbf{70\%}} 
					& $|$
					& \multicolumn{2}{c}{\textbf{80\%}} \\ 
					
					\midrule
					
					\textbf{Dataset} 
					& $|$
					& \textbf{Acc} & \textbf{F1} 
					& $|$
					& \textbf{Acc} & \textbf{F1} 
					& $|$
					& \textbf{Acc} & \textbf{F1} 
					& $|$
					& \textbf{Acc} & \textbf{F1} 
					& $|$
					& \textbf{Acc} & \textbf{F1} 
					& $|$
					& \textbf{Acc} & \textbf{F1} 
					& $|$
					& \textbf{Acc} & \textbf{F1} 
					& $|$
					& \textbf{Acc} & \textbf{F1} \\ 
					
					\midrule
					
					CICIDS2017 \cite{CICIDS2017} 
					& $|$
					& 97.12 & 97.27
					& $|$
					& 97.41 & 97.46
					& $|$
					& 97.32 & 97.38
					& $|$
					& 97.06 & 96.98
					& $|$
					& 96.09 & 96.09
					& $|$
					& 95.77 & 95.94
					& $|$
					& 93.41 & 94.33
					& $|$
					& 89.23 & 90.43 \\ 
					
					CICMalDroid-2020 \cite{CICMalDroid-2020} 
					& $|$
					& 89.09 & 89.06
					& $|$
					& 86.37 & 86.35
					& $|$
					& 84.35 & 84.32
					& $|$
					& 82.93 & 82.84
					& $|$
					& 76.20 & 75.89
					& $|$
					& 63.79 & 63.43
					& $|$
					& 52.88 & 52.29
					& $|$
					& 16.20 & 14.84 \\ 
					
					NSL-KDD \cite{NSL-KDD} 
					& $|$
					& 97.08 & 98.11
					& $|$
					& 97.26 & 98.08
					& $|$
					& 97.25 & 98.03
					& $|$
					& 97.08 & 97.95
					& $|$
					& 93.99 & 95.79
					& $|$
					& 94.92 & 96.10
					& $|$
					& 89.09 & 92.13
					& $|$
					& 13.24 & 16.95 \\ 
					
					DDoS2019-binary \cite{DDoS2019} 
					& $|$
					& 99.15 & 99.15
					& $|$
					& 97.90 & 97.89
					& $|$
					& 95.17 & 95.16
					& $|$
					& 87.82 & 87.73
					& $|$
					& 43.75 & 43.71
					& $|$
					& 13.15 & 12.02
					& $|$
					& 5.50 & 5.30
					& $|$
					& 2.32 & 2.30 \\ 
					
					\bottomrule
				\end{tabular} 
			} 
		\end{threeparttable} 
	\end{center} 
	\vspace{-0.6cm} 
\end{table*}


\begin{table*}[ht] 
	\footnotesize
	\begin{center} 
		\caption{Classification performance of Argus on different datasets under asymmetric noise scenarios.} 
		\label{closed_world_asymmetric} 
		\setlength{\tabcolsep}{5pt} 
		\renewcommand{\arraystretch}{1.2} 
		
		\begin{threeparttable} 
			\resizebox{\textwidth}{!}{ 
				\begin{tabular}{lc cc c cc c cc c cc c cc c cc c cc c cc} 
					\toprule
					
					\textbf{Noise Type} 
					& $|$
					& \multicolumn{23}{c}{\textbf{Asymmetric}} \\ 
					
					\midrule
					
					\textbf{Noise Ratio} 
					& $|$
					& \multicolumn{2}{c}{\textbf{10\%}} 
					& $|$
					& \multicolumn{2}{c}{\textbf{20\%}} 
					& $|$
					& \multicolumn{2}{c}{\textbf{30\%}} 
					& $|$
					& \multicolumn{2}{c}{\textbf{40\%}} 
					& $|$
					& \multicolumn{2}{c}{\textbf{50\%}} 
					& $|$
					& \multicolumn{2}{c}{\textbf{60\%}} 
					& $|$
					& \multicolumn{2}{c}{\textbf{70\%}} 
					& $|$
					& \multicolumn{2}{c}{\textbf{80\%}} \\ 
					
					\midrule
					
					\textbf{Dataset} 
					& $|$
					& \textbf{Acc} & \textbf{F1} 
					& $|$
					& \textbf{Acc} & \textbf{F1} 
					& $|$
					& \textbf{Acc} & \textbf{F1} 
					& $|$
					& \textbf{Acc} & \textbf{F1} 
					& $|$
					& \textbf{Acc} & \textbf{F1} 
					& $|$
					& \textbf{Acc} & \textbf{F1} 
					& $|$
					& \textbf{Acc} & \textbf{F1} 
					& $|$
					& \textbf{Acc} & \textbf{F1} \\ 
					
					\midrule
					
					CICIDS2017 \cite{CICIDS2017} 
					& $|$
					& 91.90 & 93.58
					& $|$
					& 90.73 & 92.63
					& $|$
					& 89.95 & 91.73
					& $|$
					& 89.35 & 91.05
					& $|$
					& 89.27 & 91.12
					& $|$
					& 84.99 & 87.15
					& $|$
					& 87.07 & 89.15
					& $|$
					& 79.94 & 82.23 \\ 
					
					CICMalDroid-2020 \cite{CICMalDroid-2020} 
					& $|$
					& 90.17 & 90.12
					& $|$
					& 90.64 & 90.58
					& $|$
					& 89.74 & 89.62
					& $|$
					& 89.87 & 89.77
					& $|$
					& 88.49 & 88.19
					& $|$
					& 86.68 & 86.29
					& $|$
					& 86.72 & 86.04
					& $|$
					& 84.74 & 83.64 \\ 
					
					NSL-KDD \cite{NSL-KDD} 
					& $|$
					& 99.37 & 99.41
					& $|$
					& 98.36 & 98.70
					& $|$
					& 95.53 & 96.69
					& $|$
					& 97.74 & 98.39
					& $|$
					& 89.07 & 92.88
					& $|$
					& 96.01 & 97.02
					& $|$
					& 98.79 & 98.91
					& $|$
					& 96.24 & 96.93 \\ 
					
					DDoS2019-binary \cite{DDoS2019} 
					& $|$
					& 99.22 & 99.22
					& $|$
					& 99.17 & 99.17
					& $|$
					& 98.92 & 98.92
					& $|$
					& 99.00 & 99.00
					& $|$
					& 98.92 & 98.92
					& $|$
					& 99.05 & 99.05
					& $|$
					& 99.02 & 99.02
					& $|$
					& 98.27 & 98.27 \\ 
					
					\bottomrule
				\end{tabular} 
			} 
		\end{threeparttable} 
	\end{center} 
	\vspace{-0.6cm} 
\end{table*}


The results in Tables~\ref{closed_world_symmetric} and~\ref{closed_world_asymmetric} demonstrate the strong generalization ability of Argus across different datasets. Under symmetric noise, Argus maintains high Acc and F1 on CICIDS2017 and NSL-KDD when the noise ratio is below 60\%, indicating its effectiveness in suppressing random label noise. However, performance drops noticeably under extremely high noise ratios, especially on DDoS2019-binary and CICMalDroid-2020, suggesting that excessive random noise can severely damage class-discriminative boundaries.

Under asymmetric noise, Argus shows more stable performance. CICMalDroid-2020 and DDoS2019-binary maintain consistently high Acc and F1 across all noise ratios, with DDoS2019-binary remaining close to 99\%. Overall, these results indicate that Argus can generalize well to datasets with different traffic distributions and task types, demonstrating strong cross-dataset robustness against noisy labels.



\section{Conclusion}\label{conclusion}
In the dynamic network environment, we



%In the dynamic network environment, network intrusion detection systems face the dual challenges of label noise and unknown attack detection. This paper proposes a novel anti-noise detection framework, AEGIS-Net. AEGIS-Net dynamically selects multiple prototypes for each class through a multi-prototype strategy to achieve accurate noise label correction. It also designs a model-agnostic k-nearest neighbors distance detection paradigm, which determines unknown attacks by calculating the relative distance ratio between the test sample and known classes, eliminating the dependence on traffic distribution assumptions common in traditional methods. This makes the model more flexible and versatile. We prove, through the feature space stability hypothesis, that the detection threshold is independent of the noise rate, mathematically verifying the anti-noise property of unknown attack detection. Additionally, a series of experiments demonstrate the effectiveness of AEGIS-Net in known attack classification, unknown attack detection, and cross-scenario generalization ability. Moreover, we show performance improvement compared to the state-of-the-art baseline methods.


%%% algorithm. 算法
%%% Please use Algorithm~\ref{alg1} in the text (``~'' is not omitted)
%%% 请在文中使用Algorithm~\ref{alg1}, 不要省略``~'', 它能够使Algorithm和算法编号不断行
%%% 双栏排版的通栏算法, 使用\begin{algorithm*}...\end{algorithm*}环境, 其他不变
%%% 双栏排版的单栏算法, 使用\begin{algorithm}[H]...\end{algorithm}环境, 其他不变


%\begin{algorithm}
%%\floatname{algorithm}{Algorithm}%更改算法前缀名称
%%\renewcommand{\algorithmicrequire}{\textbf{Input:}}%更改输入名称
%%\renewcommand{\algorithmicensure}{\textbf{Output:}}%更改输出名称
%\footnotesize
%\caption{Algorithm caption}
%\label{alg1}
%\begin{algorithmic}[1]
%    \REQUIRE $n \geq 0 \vee x \neq 0$;
%    \ENSURE $y = x^n$;
%    \STATE $y \Leftarrow 1$;
%    \IF{$n < 0$}
%        \STATE $X \Leftarrow 1 / x$;
%        \STATE $N \Leftarrow -n$;
%    \ELSE
%        \STATE $X \Leftarrow x$;
%        \STATE $N \Leftarrow n$;
%    \ENDIF
%    \WHILE{$N \neq 0$}
%        \IF{$N$ is even}
%            \STATE $X \Leftarrow X \times X$;
%            \STATE $N \Leftarrow N / 2$;
%        \ELSE[$N$ is odd]
%            \STATE $y \Leftarrow y \times X$;
%            \STATE $N \Leftarrow N - 1$;
%        \ENDIF
%    \ENDWHILE
%\end{algorithmic}
%\end{algorithm}


% \newpage

%\section{More information}
%The examples at the bottom of the .tex file can help you when preparing your manuscript. We are appreciate your effort to follow our style~\cite{1,2}.


%\theorem[name]{It is a test theorem.}

%\begin{table}[!t]
%\footnotesize
%\caption{Tabel caption}
%\label{tab1}
%\tabcolsep 49pt %space between two columns. 用于调整列间距
%\begin{tabular*}{\textwidth}{cccc}
%\toprule
%  Title a & Title b & Title c & Title d \\\hline
%  Aaa & Bbb & Ccc\footnote & Ddd\\
%  Aaa & Bbb & Ccc\footnote & Ddd\\
%  Aaa & Bbb & Ccc & Ddd\\
%\bottomrule
%\end{tabular*}
%\end{table}



%\footnotetext[1]{Footnote 1.}
%\footnotetext[2]{Footnote 2.}

%%%%%%%%%%%%%%%%%%%%%%%%%%%%%%%%%%%%%%%%%%%%%%%%%%%%%%%
%%% Acknowledgements. 致谢
%%%%%%%%%%%%%%%%%%%%%%%%%%%%%%%%%%%%%%%%%%%%%%%%%%%%%%%
\Acknowledgements{This work was supported by the National Natural Science Foundation of China (Grant Nos. 00000000 and 11111111).}


%%%%%%%%%%%%%%%%%%%%%%%%%%%%%%%%%%%%%%%%%%%%%%%%%%%%%%%
%%% Supplements. 补充材料, 非必选
%%%%%%%%%%%%%%%%%%%%%%%%%%%%%%%%%%%%%%%%%%%%%%%%%%%%%%%
\Supplements{Appendix A.}

%%%%%%%%%%%%%%%%%%%%%%%%%%%%%%%%%%%%%%%%%%%%%%%%%%%%%%%
%%% Open Access Funding Note. 开放获取基金来源说明
%%%%%%%%%%%%%%%%%%%%%%%%%%%%%%%%%%%%%%%%%%%%%%%%%%%%%%%
%\FundingNote{\textbf{Open Access} funding enabled and organized by CAUL and its Member Institutions.}

%%%%%%%%%%%%%%%%%%%%%%%%%%%%%%%%%%%%%%%%%%%%%%%%%%%%%%%
%%% Reference section. 参考文献
%%% Citation in the content using "some words~\cite{1,2}".
%%% ~ is needed to make the reference number is on the same line with the word before it.
%%% Using scis.bst to format the style if the ref.bib file is included, e.g.
\bibliographystyle{scis}
\bibliography{ref}
%%%%%%%%%%%%%%%%%%%%%%%%%%%%%%%%%%%%%%%%%%%%%%%%%%%%%%%
%\begin{thebibliography}{99}
%
%\bibitem{1} Author A, Author B, Author C. Reference title. Journal, 2024, 38: 13--28
%
%\bibitem{2} Author A, Author B, Author C, et al. Reference title. In: Proceedings of Conference, Place, 2024. 6--12
%
%\end{thebibliography}

%%%%%%%%%%%%%%%%%%%%%%%%%%%%%%%%%%%%%%%%%%%%%%%%%%%%%%%
%%% Appendix sections. 附录章节, 非必选
%%%%%%%%%%%%%%%%%%%%%%%%%%%%%%%%%%%%%%%%%%%%%%%%%%%%%%%
%\begin{appendix}
%\section{Name}

%\end{appendix}

\end{document}


%%%%%%%%%%%%%%%%%%%%%%%%%%%%%%%%%%%%%%%%%%%%%%%%%%%%%%%
%%% Some latex examples for this template
%%% 本模板使用的latex排版示例
%%%%%%%%%%%%%%%%%%%%%%%%%%%%%%%%%%%%%%%%%%%%%%%%%%%%%%%

%%% sections. 章节
\section{Section title}
\subsection{Subsection title}
\subsubsection{Subsubsection title}


%%% list. 普通列表
\begin{itemize}
\item Aaa aaa.
\item Bbb bbb.
\item Ccc ccc.
\end{itemize}

%%% list-with-other-label. 自编号列表
\begin{itemize}
\itemindent 2.8em
\item[(1)] Aaa aaa.
\item[(2)] Bbb bbb.
\item[(3)] Ccc ccc.
\end{itemize}

%%% Theorem-like section. 定义、定理、引理、推论
%%% [name] can be ommited. [name]可以省略
\assumption[name]{Content.}
\corollary[name]{Content.}
\definition[name]{Content.}
\example[name]{Content.}
\lemma[name]{Content.}
\problem[name]{Content.}
\proposition[name]{Content.}
\remark[name]{Content.}
\theorem[name]{Content.}

%%% If other type of Theorem-like section, 若使用定理样式的其他前缀
%%% please use \newtheorem command before \begin{document}. 在 "作者附加的定义" 处加入\newtheorem命令
%%% For example, "Definition" section is commanded by: 例如 "Definition" 是由以下命令定义的
\newtheorem{definition}{Definition}


%%% 1-figure. 单图
%%% Please use Figure~\ref{fig1} in the text (``~'' is not omitted)
%%% 请在文中使用Figure~\ref{fig1}, 不要省略``~'', 它能够使Figure和图号不断行
%%% 双栏排版的通栏图片, 使用\begin{figure*}...\end{figure*}环境, 其他不变
%%% 双栏排版的单栏图片, 使用\begin{figure}[H]...\end{figure}环境, 其他不变
\begin{figure}[!t]
\centering
\includegraphics{fig1.eps}
\caption{Caption 1.}
\label{fig1}
\end{figure}

%%% 2-figures-in-1-row. 并排图
%%% Please use Figure~\ref{fig1} in the text (``~'' is not omitted)
%%% 请在文中使用Figure~\ref{fig1}, 不要省略``~'', 它能够使Figure和图号不断行
\begin{figure}[!t]
\centering
\begin{minipage}[c]{0.48\textwidth}
\centering
\includegraphics{fig1.eps}
\end{minipage}
\hspace{0.02\textwidth}
\begin{minipage}[c]{0.48\textwidth}
\centering
\includegraphics{fig2.eps}
\end{minipage}\\[3mm]
\begin{minipage}[t]{0.48\textwidth}
\centering
\caption{Caption 1.}
\label{fig1}
\end{minipage}
\hspace{0.02\textwidth}
\begin{minipage}[t]{0.48\textwidth}
\centering
\caption{Caption 2.}
\label{fig2}
\end{minipage}
\end{figure}

%%% 2-subfigures-in-1-row. 并排子图
%%% Please use Figure~\ref{fig1} in the text (``~'' is not omitted)
%%% 请在文中使用Figure~\ref{fig1}, 不要省略``~'', 它能够使Figure和图号不断行
\begin{figure}[!t]
\centering
\begin{minipage}[c]{0.48\textwidth}
\centering
\includegraphics{subfig1.eps}
\end{minipage}
\hspace{0.02\textwidth}
\begin{minipage}[c]{0.48\textwidth}
\centering
\includegraphics{subfig2.eps}
\end{minipage}
\caption{Caption 1. (a) Subfig1 caption; (b) subfig2 caption.}
\label{fig1}
\end{figure}


%%% algorithm. 算法
%%% Please use Algorithm~\ref{alg1} in the text (``~'' is not omitted)
%%% 请在文中使用Algorithm~\ref{alg1}, 不要省略``~'', 它能够使Algorithm和算法编号不断行
%%% 双栏排版的通栏算法, 使用\begin{algorithm*}...\end{algorithm*}环境, 其他不变
%%% 双栏排版的单栏算法, 使用\begin{algorithm}[H]...\end{algorithm}环境, 其他不变
\begin{algorithm}
%\floatname{algorithm}{Algorithm}%更改算法前缀名称
%\renewcommand{\algorithmicrequire}{\textbf{Input:}}%更改输入名称
%\renewcommand{\algorithmicensure}{\textbf{Output:}}%更改输出名称
\footnotesize
\caption{Algorithm caption}
\label{alg1}
\begin{algorithmic}[1]
    \REQUIRE $n \geq 0 \vee x \neq 0$;
    \ENSURE $y = x^n$;
    \STATE $y \Leftarrow 1$;
    \IF{$n < 0$}
        \STATE $X \Leftarrow 1 / x$;
        \STATE $N \Leftarrow -n$;
    \ELSE
        \STATE $X \Leftarrow x$;
        \STATE $N \Leftarrow n$;
    \ENDIF
    \WHILE{$N \neq 0$}
        \IF{$N$ is even}
            \STATE $X \Leftarrow X \times X$;
            \STATE $N \Leftarrow N / 2$;
        \ELSE[$N$ is odd]
            \STATE $y \Leftarrow y \times X$;
            \STATE $N \Leftarrow N - 1$;
        \ENDIF
    \ENDWHILE
\end{algorithmic}
\end{algorithm}

%%% simple-table. 简单表格
%%% Please use Table~\ref{tab1} in the text (``~'' is not omitted)
%%% 请在文中使用Table~\ref{tab1}, 不要省略``~'', 它能够使Table和表号不断行
%%% 双栏排版的通栏表格, 使用\begin{table*}...\end{table*}环境, 其他不变
%%% 双栏排版的单栏表格, 使用\begin{table}[H]...\end{table}环境, 其他不变
\begin{table}[!t]
\footnotesize
\caption{Tabel caption}
\label{tab1}
\tabcolsep 49pt %space between two columns. 用于调整列间距
\begin{tabular*}{\textwidth}{cccc}
\toprule
  Title a & Title b & Title c & Title d \\\hline
  Aaa & Bbb & Ccc & Ddd\\
  Aaa & Bbb & Ccc & Ddd\\
  Aaa & Bbb & Ccc & Ddd\\
\bottomrule
\end{tabular*}
\end{table}

%%% linefeed-table. 换行表格
%%% Please use Table~\ref{tab1} in the text (``~'' is not omitted)
%%% 请在文中使用Table~\ref{tab1}, 不要省略``~'', 它能够使Table和表号不断行
\begin{table}[!t]
\footnotesize
\caption{Tabel caption}
\label{tab1}
\def\tabblank{\hspace*{10mm}} %blank leaving of both side of the table. 左右两边的留白
\begin{tabularx}{\textwidth} %using p{?mm} to define the width of a column. 用p{?mm}控制列宽
{@{\tabblank}@{\extracolsep{\fill}}cccp{100mm}@{\tabblank}}
\toprule
  Title a & Title b & Title c & Title d \\\hline
  Aaa & Bbb & Ccc & Ddd ddd ddd ddd.

  Ddd ddd ddd ddd ddd ddd ddd ddd ddd ddd ddd ddd ddd ddd ddd ddd ddd ddd ddd ddd ddd ddd ddd ddd ddd ddd ddd ddd ddd ddd ddd.\\
  Aaa & Bbb & Ccc & Ddd ddd ddd ddd.\\
  Aaa & Bbb & Ccc & Ddd ddd ddd ddd.\\
\bottomrule
\end{tabularx}
\end{table}

%%% 1-line-equation. 单行公式
%%% Please use (\ref{eq1}) in the text. Eq.~(\ref{eq1}) is used when the equation number is the beginning of the sentence (``~'' is not omitted)
%%% 请在文中使用(\ref{eq1}), 如果是句子开头, 使用Eq.~(\ref{eq1}), 不要省略``~'', 它能够使Eq.和公式号不断行
\begin{equation}
A(d,f)=d^{l}a^{d}(f),
\label{eq1}
\end{equation}

%%% 1-line-equation-without-count-number. 不编号单行公式
\begin{equation}
\nonumber
A(d,f)=d^{l}a^{d}(f),
\end{equation}

%%% equation-array. 公式组
\begin{eqnarray}
\nonumber
&X=[x_{11},x_{12},\ldots,x_{ij},\ldots ,x_{n-1,n}]^{\rm T},\\
\nonumber
&\varepsilon=[e_{11},e_{12},\ldots ,e_{ij},\ldots ,e_{n-1,n}],\\
\nonumber
&T=[t_{11},t_{12},\ldots ,t_{ij},\ldots ,t_{n-1,n}].
\end{eqnarray}

%%% conditional-equation. 条件公式
%%% Please use (\ref{eq1}) in the text. Eq.~(\ref{eq1}) is used when the equation number is the beginning of the sentence (``~'' is not omitted)
%%% 请在文中使用(\ref{eq1}), 如果是句子开头, 使用Eq.~(\ref{eq1}), 不要省略``~'', 它能够使Eq.和公式号不断行
\begin{eqnarray}
\sum_{j=1}^{n}x_{ij}-\sum_{k=1}^{n}x_{ki}=
\left\{
\begin{aligned}
1,&\quad i=1,\\
0,&\quad i=2,\ldots ,n-1,\\
-1,&\quad i=n.
\end{aligned}
\right.
\label{eq1}
\end{eqnarray}

%%% Other. 其他格式
\proof %proof. 证明
\footnote{Comments.} %footnote. 脚注
\raisebox{-1pt}[0mm][0mm]{xxxx} %put xxxx upper or lower. 控制xxxx的垂直位置
